\documentclass[11pt,a4paper]{article}

%====================================================
% PACKAGES
%====================================================
\usepackage[utf8]{inputenc}
\usepackage[T1]{fontenc}
\usepackage[french]{babel}
\usepackage[a4paper,margin=1.7cm]{geometry}
\usepackage{amsmath,amssymb,amsthm,mathtools}
\usepackage{physics}
\usepackage{siunitx}
\usepackage{lmodern}
\usepackage{xcolor}
\usepackage{hyperref}
\usepackage{fancyhdr}
\usepackage{lastpage}
\usepackage{enumitem}
\usepackage{array}
\usepackage{tabularx}
\usepackage{multicol}
\usepackage{tikz}
\usepackage{pgfplots}
\pgfplotsset{compat=1.18}
\sisetup{locale=FR}

%====================================================
% MISE EN PAGE
%====================================================
\hypersetup{
    colorlinks=true,
    linkcolor=blue!50!black,
    urlcolor=blue!50!black
}
\setlength{\headheight}{14pt}
\pagestyle{fancy}
\fancyhf{}
\fancyhead[L]{Terminale générale -- meca--pgb}
\fancyhead[R]{Page \thepage/\pageref{LastPage}}
\fancyfoot[C]{Progression complète avec corrigé final}

%====================================================
% ENVIRONNEMENTS
%====================================================
\theoremstyle{definition}
\newtheorem{exercice}{Exercice}[section]

%====================================================
% COMMANDES
%====================================================
\newcommand{\etoiles}[1]{\hfill{\small\textbf{#1}}}
\newcommand{\vF}{\textbf{Vrai/Faux}}
\newcommand{\QCM}{\textbf{QCM}}
\newcommand{\R}{\mathbb{R}}
\newcommand{\vect}[1]{\overrightarrow{#1}}
\newcommand{\g}{\vec{g}}

\begin{document}

\begin{center}
    {\LARGE \textbf{TD complet -- Mécanique}}\\[0.2cm]
    {\large Terminale spécialité Physique-Chimie}\\[0.2cm]
    {\small Mouvement et interactions, lois de Newton, chute libre, projectiles, quantité de mouvement, mouvement circulaire}\\[0.2cm]
    {\small Toutes les corrections sont regroupées à la fin.}
\end{center}

\vspace{0.3cm}

\begin{center}
\renewcommand{\arraystretch}{1.3}
\begin{tabular}{|c|l|}
\hline
$\star$ & application directe / très accessible \\
\hline
$\star\star$ & classique Terminale \\
\hline
$\star\star\star$ & bon niveau Terminale / type bac standard \\
\hline
$\star\star\star\star$ & difficile / demande du recul \\
\hline
$\star\star\star\star\star$ & approfondissement très poussé \\
\hline
\end{tabular}
\end{center}

\tableofcontents
\newpage

%====================================================
\section{10 Vrai / Faux}
%====================================================

\begin{exercice}[Vrai/Faux 1 \etoiles{$\star$}]
Pour chaque affirmation, dire si elle est vraie ou fausse et justifier brièvement.
\begin{enumerate}[label=\alph*)]
    \item Si la somme des forces extérieures est nulle, alors le système est immobile.
    \item Un mouvement rectiligne uniforme est associé à une accélération nulle.
    \item Le poids d’un corps est proportionnel à sa masse.
    \item Le poids s’exprime en \si{kg}.
\end{enumerate}
\end{exercice}

\begin{exercice}[Vrai/Faux 2 \etoiles{$\star$}]
\begin{enumerate}[label=\alph*)]
    \item Le vecteur vitesse est tangent à la trajectoire.
    \item Le vecteur accélération est toujours parallèle au vecteur vitesse.
    \item Dans un mouvement circulaire uniforme, la vitesse est constante.
    \item Dans un mouvement circulaire uniforme, l’accélération est nulle.
\end{enumerate}
\end{exercice}

\begin{exercice}[Vrai/Faux 3 \etoiles{$\star\star$}]
\begin{enumerate}[label=\alph*)]
    \item La deuxième loi de Newton s’écrit $\sum \vec{F}=m\vec{a}$ pour un système de masse constante.
    \item La troisième loi de Newton relie deux forces qui s’exercent sur le même objet.
    \item La réaction du support peut compenser le poids.
    \item Un projectile sans frottements a une trajectoire parabolique.
\end{enumerate}
\end{exercice}

\begin{exercice}[Vrai/Faux 4 \etoiles{$\star\star$}]
\begin{enumerate}[label=\alph*)]
    \item En chute libre sans frottements, seule la force poids agit sur le système.
    \item En lancer horizontal, la vitesse horizontale reste constante si l’air est négligé.
    \item En lancer horizontal, l’accélération horizontale vaut $g$.
    \item En chute libre verticale, la vitesse varie linéairement avec le temps.
\end{enumerate}
\end{exercice}

\begin{exercice}[Vrai/Faux 5 \etoiles{$\star\star$}]
\begin{enumerate}[label=\alph*)]
    \item Une force peut modifier la direction du mouvement sans modifier la valeur de la vitesse.
    \item Si un système avance, alors une force agit forcément dans le sens du mouvement.
    \item Si la vitesse d’un système augmente, alors son accélération est forcément constante.
    \item Un objet immobile subit toujours une résultante des forces nulle.
\end{enumerate}
\end{exercice}

\begin{exercice}[Vrai/Faux 6 \etoiles{$\star\star\star$}]
\begin{enumerate}[label=\alph*)]
    \item Dans un référentiel terrestre usuel, on peut souvent assimiler le référentiel à un référentiel galiléen.
    \item Si la résultante des forces est constante, alors l’accélération est constante.
    \item Si l’accélération est constante, la trajectoire est forcément rectiligne.
    \item Dans un champ électrique uniforme, une charge négative est accélérée dans le sens opposé au champ.
\end{enumerate}
\end{exercice}

\begin{exercice}[Vrai/Faux 7 \etoiles{$\star\star\star$}]
\begin{enumerate}[label=\alph*)]
    \item La quantité de mouvement s’écrit $\vec{p}=m\vec{v}$.
    \item Si la vitesse double, la quantité de mouvement double.
    \item La quantité de mouvement s’exprime en newton.
    \item Pour une masse constante, $\sum \vec{F}=\dv{\vec{p}}{t}$.
\end{enumerate}
\end{exercice}

\begin{exercice}[Vrai/Faux 8 \etoiles{$\star\star\star$}]
\begin{enumerate}[label=\alph*)]
    \item Dans un virage à vitesse constante, l’accélération est dirigée vers le centre.
    \item L’accélération centripète vaut $\dfrac{v^2}{R}$.
    \item Si le rayon du virage augmente à vitesse fixée, l’accélération centripète augmente.
    \item La force résultante responsable d’un mouvement circulaire uniforme est tangentielle.
\end{enumerate}
\end{exercice}

\begin{exercice}[Vrai/Faux 9 \etoiles{$\star\star\star$}]
\begin{enumerate}[label=\alph*)]
    \item Si un objet glisse à vitesse constante sur une table horizontale, alors la somme des forces peut être nulle.
    \item La masse dépend du lieu où se trouve l’objet.
    \item Le poids dépend du lieu où se trouve l’objet.
    \item Sur la Lune, un objet a la même masse que sur Terre.
\end{enumerate}
\end{exercice}

\begin{exercice}[Vrai/Faux 10 \etoiles{$\star\star\star\star$}]
\begin{enumerate}[label=\alph*)]
    \item Deux forces d’action-réaction ont même intensité, même direction, sens opposés.
    \item Deux forces d’action-réaction peuvent s’annuler.
    \item Un système peut avoir une vitesse non nulle et une accélération nulle.
    \item Un système peut avoir une vitesse nulle et une accélération non nulle.
\end{enumerate}
\end{exercice}

%====================================================
\section{10 QCM}
%====================================================

\begin{exercice}[QCM 1 \etoiles{$\star$}]
L’unité de l’accélération est :
\begin{multicols}{4}
\begin{enumerate}[label=\Alph*.]
    \item \si{N}
    \item \si{m.s^{-1}}
    \item \si{m.s^{-2}}
    \item \si{kg.m.s^{-1}}
\end{enumerate}
\end{multicols}
\end{exercice}

\begin{exercice}[QCM 2 \etoiles{$\star$}]
Le poids d’un corps de masse $m$ vaut :
\begin{multicols}{4}
\begin{enumerate}[label=\Alph*.]
    \item $\vec{P}=m\vec{a}$
    \item $\vec{P}=m\vec{g}$
    \item $\vec{P}=g\vec{v}$
    \item $P=mg^2$
\end{enumerate}
\end{multicols}
\end{exercice}

\begin{exercice}[QCM 3 \etoiles{$\star\star$}]
Dans un mouvement rectiligne uniforme :
\begin{multicols}{2}
\begin{enumerate}[label=\Alph*.]
    \item la vitesse varie
    \item l’accélération est constante non nulle
    \item l’accélération est nulle
    \item la trajectoire est circulaire
\end{enumerate}
\end{multicols}
\end{exercice}

\begin{exercice}[QCM 4 \etoiles{$\star\star$}]
La deuxième loi de Newton s’applique, dans ce chapitre, dans :
\begin{multicols}{2}
\begin{enumerate}[label=\Alph*.]
    \item tout référentiel quelconque sans précaution
    \item un référentiel galiléen
    \item uniquement le référentiel héliocentrique
    \item uniquement pour les objets immobiles
\end{enumerate}
\end{multicols}
\end{exercice}

\begin{exercice}[QCM 5 \etoiles{$\star\star$}]
Pour un lancer horizontal sans frottements :
\begin{multicols}{2}
\begin{enumerate}[label=\Alph*.]
    \item $a_x=g$
    \item $a_z=0$
    \item $v_x$ est constante
    \item la trajectoire est une droite
\end{enumerate}
\end{multicols}
\end{exercice}

\begin{exercice}[QCM 6 \etoiles{$\star\star\star$}]
Dans un mouvement circulaire uniforme :
\begin{multicols}{2}
\begin{enumerate}[label=\Alph*.]
    \item la vitesse vectorielle est constante
    \item l’accélération est tangentielle
    \item l’accélération est centripète
    \item la résultante des forces est nulle
\end{enumerate}
\end{multicols}
\end{exercice}

\begin{exercice}[QCM 7 \etoiles{$\star\star\star$}]
La quantité de mouvement d’un système est :
\begin{multicols}{2}
\begin{enumerate}[label=\Alph*.]
    \item $\vec{p}=m\vec{a}$
    \item $\vec{p}=m\vec{v}$
    \item $\vec{p}=q\vec{E}$
    \item $\vec{p}=m\vec{g}$
\end{enumerate}
\end{multicols}
\end{exercice}

\begin{exercice}[QCM 8 \etoiles{$\star\star\star$}]
Une charge négative placée dans un champ électrique uniforme $\vec{E}$ :
\begin{multicols}{2}
\begin{enumerate}[label=\Alph*.]
    \item ne subit aucune force
    \item est accélérée dans le sens de $\vec{E}$
    \item est accélérée en sens opposé à $\vec{E}$
    \item reste immobile
\end{enumerate}
\end{multicols}
\end{exercice}

\begin{exercice}[QCM 9 \etoiles{$\star\star\star$}]
Si $\sum \vec{F}=\vec{0}$, alors :
\begin{multicols}{2}
\begin{enumerate}[label=\Alph*.]
    \item le système est au repos
    \item le système a forcément une vitesse constante
    \item le système a une accélération nulle
    \item le système ralentit
\end{enumerate}
\end{multicols}
\end{exercice}

\begin{exercice}[QCM 10 \etoiles{$\star\star\star\star$}]
Un objet lancé verticalement vers le haut, sans frottements :
\begin{multicols}{2}
\begin{enumerate}[label=\Alph*.]
    \item a une accélération nulle au sommet
    \item a une vitesse nulle au sommet
    \item a encore son poids au sommet
    \item change de sens de mouvement au sommet
\end{enumerate}
\end{multicols}
\end{exercice}

%====================================================
\section{Grands exercices type bac}
%====================================================

\begin{exercice}[Type bac 1 -- Plongeon de haut vol \etoiles{$\star\star\star$}]
Un sportif effectue un plongeon depuis une plateforme située à \SI{28}{m} au-dessus de l’eau.
À la sortie du plongeoir, sa vitesse initiale est de norme $v_0=\SI{2,6}{m.s^{-1}}$ et fait un angle
$\alpha=35^\circ$ au-dessus de l’horizontale.

On travaille dans le référentiel terrestre supposé galiléen. On note $g=\SI{9,81}{m.s^{-2}}$.
On néglige dans un premier temps l’action de l’air.

On choisit un repère $(O;\vec{i},\vec{j})$ où $O$ est la position initiale du plongeur, $\vec{i}$ horizontal
et $\vec{j}$ vertical vers le haut.

\begin{enumerate}[label=\arabic*.]
    \item Décomposer le vecteur vitesse initiale selon les axes du repère.
    \item Faire le bilan des forces exercées sur le plongeur pendant la phase aérienne.
    \item Justifier que le mouvement peut être modélisé comme une chute libre.
    \item Écrire la deuxième loi de Newton puis déterminer les coordonnées du vecteur accélération.
    \item Établir les expressions de $v_x(t)$ et $v_y(t)$.
    \item En déduire les équations horaires $x(t)$ et $y(t)$.
    \item Établir l’équation cartésienne de la trajectoire.
    \item Montrer que la trajectoire est parabolique.
    \item Déterminer la date d’entrée dans l’eau.
    \item Déterminer les composantes de la vitesse à l’entrée dans l’eau.
    \item En déduire la norme de la vitesse d’impact.
    \item Le plongeur affirme que sa vitesse d’entrée dans l’eau est d’environ \SI{24}{m.s^{-1}}.
    Discuter cette affirmation.
    \item On mesure expérimentalement une durée de chute $\Delta t_{\mathrm{exp}}=\SI{2,8}{s}$ avec une
    incertitude de \SI{0,3}{s}. Le modèle précédent est-il compatible avec la mesure ?
    \item Citer au moins deux causes possibles d’écart entre le modèle et la réalité.
\end{enumerate}
\end{exercice}

\begin{exercice}[Type bac 2 -- Service au volley-ball \etoiles{$\star\star\star$}]
Lors d’un service smashé, un joueur frappe un ballon depuis une hauteur de \SI{2,2}{m}.
Au départ, la vitesse du centre du ballon a pour norme $v_0=\SI{18}{m.s^{-1}}$ et fait un angle
$\alpha=12^\circ$ au-dessus de l’horizontale.

Le filet est situé à \SI{9,0}{m} du joueur et sa hauteur est de \SI{2,43}{m}.
Le terrain adverse mesure encore \SI{9,0}{m} après le filet.

On néglige les frottements de l’air.

\begin{enumerate}[label=\arabic*.]
    \item Décomposer la vitesse initiale.
    \item Établir les équations horaires du mouvement.
    \item Établir l’équation de la trajectoire.
    \item Calculer la hauteur du ballon lorsqu’il passe à l’aplomb du filet.
    \item Le ballon franchit-il le filet ?
    \item Déterminer la position horizontale du point de chute.
    \item Le service tombe-t-il dans le terrain adverse ?
    \item Déterminer la durée totale du vol.
    \item Déterminer les composantes du vecteur vitesse juste avant le contact au sol.
    \item Comment varie qualitativement la portée si :
    \begin{enumerate}[label=\alph*)]
        \item la vitesse initiale augmente ;
        \item l’angle de frappe augmente légèrement ;
        \item la hauteur initiale diminue ?
    \end{enumerate}
    \item Expliquer pourquoi le modèle de projectile sans frottement reste imparfait pour un ballon réel.
\end{enumerate}
\end{exercice}

\begin{exercice}[Type bac 3 -- Freinage d’urgence d’une voiture \etoiles{$\star\star\star\star$}]
Une voiture de masse \SI{1350}{kg} roule en ligne droite à la vitesse de \SI{90}{km.h^{-1}}.
Le conducteur perçoit un obstacle, met \SI{0,80}{s} à réagir, puis freine.
Pendant la phase de freinage, l’accélération est supposée constante et vaut
$a=\SI{-6,5}{m.s^{-2}}$.

\begin{enumerate}[label=\arabic*.]
    \item Convertir la vitesse initiale en \si{m.s^{-1}}.
    \item Calculer la distance parcourue pendant le temps de réaction.
    \item Déterminer la durée de freinage.
    \item Déterminer la distance de freinage.
    \item En déduire la distance totale d’arrêt.
    \item Déterminer la variation de quantité de mouvement de la voiture pendant le freinage.
    \item Calculer la résultante des forces extérieures pendant le freinage.
    \item Cette résultante est-elle orientée dans le sens du mouvement ou en sens opposé ?
    \item On répète l’expérience à \SI{110}{km.h^{-1}} avec le même temps de réaction et la même
    décélération. Refaire les calculs de distance d’arrêt.
    \item Comparer les distances d’arrêt et commenter le danger de l’augmentation de vitesse.
    \item Expliquer pourquoi la distance de freinage n’est pas proportionnelle à la vitesse initiale.
\end{enumerate}
\end{exercice}

\begin{exercice}[Type bac 4 -- Virage sur route et adhérence \etoiles{$\star\star\star\star$}]
Une voiture de masse \SI{1000}{kg} aborde un virage circulaire horizontal de rayon \SI{45}{m}
à la vitesse constante de \SI{54}{km.h^{-1}}.

On suppose que la force latérale responsable du virage est fournie par l’adhérence des pneus sur la route.

\begin{enumerate}[label=\arabic*.]
    \item Convertir la vitesse en \si{m.s^{-1}}.
    \item Déterminer l’accélération centripète.
    \item En déduire la résultante des forces exercées sur la voiture.
    \item Préciser la direction et le sens de cette résultante.
    \item Pourquoi la voiture peut-elle avoir une vitesse de norme constante tout en étant accélérée ?
    \item Si la vitesse double, par combien est multipliée l’accélération centripète ?
    \item Même question pour la force latérale nécessaire.
    \item On considère maintenant un second virage de rayon \SI{90}{m}, parcouru à la même vitesse.
    Comparer les accélérations centripètes dans les deux cas.
    \item Expliquer physiquement pourquoi les virages serrés sont plus difficiles à négocier.
    \item Discuter le rôle de l’adhérence et les limites du modèle en cas de pluie.
\end{enumerate}
\end{exercice}

\begin{exercice}[Type bac 5 -- Fusée au décollage \etoiles{$\star\star\star\star$}]
On modélise le début du décollage d’une petite fusée par un mouvement vertical.
La masse de la fusée est supposée constante pendant les premières secondes et vaut \SI{800}{kg}.
La poussée exercée par les gaz vaut \SI{12000}{N}, verticale vers le haut.
On néglige les frottements de l’air.

\begin{enumerate}[label=\arabic*.]
    \item Faire le bilan des forces exercées sur la fusée.
    \item Calculer son poids.
    \item Déterminer la résultante des forces.
    \item En déduire l’accélération de la fusée.
    \item La fusée peut-elle décoller si la poussée vaut seulement \SI{7000}{N} ? Justifier.
    \item En supposant l’accélération constante, déterminer la vitesse au bout de \SI{4,0}{s}.
    \item Déterminer l’altitude gagnée pendant cette durée.
    \item Comment évoluerait qualitativement le mouvement si l’on prenait en compte les frottements ?
    \item Expliquer pourquoi, dans un modèle plus réaliste, la masse de la fusée ne peut pas être supposée constante.
\end{enumerate}
\end{exercice}

\begin{exercice}[Type bac 6 -- Particule chargée dans un champ électrique \etoiles{$\star\star\star\star$}]
Une particule de masse $m=\SI{2,0e-4}{kg}$ et de charge $q=-\SI{5,0e-6}{C}$
entre avec une vitesse initiale horizontale dans une zone où règne un champ électrique uniforme vertical
de norme $E=\SI{800}{V.m^{-1}}$.

On néglige le poids devant la force électrique dans un premier temps.

\begin{enumerate}[label=\arabic*.]
    \item Déterminer le sens de la force électrique subie par la particule.
    \item Calculer la norme de cette force.
    \item Déterminer la norme de l’accélération.
    \item Quelle analogie peut-on faire avec l’étude d’un projectile dans le champ de pesanteur ?
    \item Si la particule entre avec une vitesse horizontale $v_0=\SI{12}{m.s^{-1}}$ dans une zone de longueur
    \SI{0,50}{m}, déterminer la durée de traversée.
    \item En déduire la déviation verticale pendant la traversée.
    \item Reprendre la question précédente en tenant compte cette fois du poids.
    \item Comparer l’importance de la force électrique et du poids.
    \item Conclure sur la pertinence de l’hypothèse initiale.
\end{enumerate}
\end{exercice}

\begin{exercice}[Type bac 7 -- Ascenseur : immobilité, montée, freinage \etoiles{$\star\star\star$}]
Une personne de masse \SI{70}{kg} se tient debout dans un ascenseur.
On note $R$ la réaction du sol de la cabine sur la personne.

\begin{enumerate}[label=\arabic*.]
    \item Lorsque l’ascenseur est immobile, calculer le poids de la personne et déterminer $R$.
    \item Même question si l’ascenseur monte à vitesse constante.
    \item Même question si l’ascenseur accélère vers le haut avec $a=\SI{1,2}{m.s^{-2}}$.
    \item Même question si l’ascenseur freine en montant avec une accélération vers le bas de
    \SI{1,2}{m.s^{-2}}.
    \item Dans quel cas la personne se sent-elle plus lourde ?
    \item Dans quel cas se sent-elle plus légère ?
    \item Expliquer pourquoi la réaction du support n’est pas toujours égale au poids.
\end{enumerate}
\end{exercice}

\begin{exercice}[Type bac 8 -- Chute comparée Terre / Lune \etoiles{$\star\star\star$}]
On laisse tomber sans vitesse initiale une balle depuis une hauteur de \SI{10}{m} :
\begin{itemize}
    \item sur Terre, avec $g_T=\SI{9,81}{m.s^{-2}}$ ;
    \item sur la Lune, avec $g_L=\SI{1,62}{m.s^{-2}}$.
\end{itemize}

\begin{enumerate}[label=\arabic*.]
    \item Établir l’expression générale du temps de chute depuis une hauteur $h$.
    \item Calculer le temps de chute sur Terre.
    \item Calculer le temps de chute sur la Lune.
    \item Déterminer la vitesse d’impact dans chaque cas.
    \item Comparer la masse de la balle sur Terre et sur la Lune.
    \item Comparer son poids sur Terre et sur la Lune.
    \item Expliquer pourquoi les temps de chute sont différents alors que la masse ne change pas.
\end{enumerate}
\end{exercice}

\begin{exercice}[Type bac 9 -- Palet sur glace puis choc contre la bande \etoiles{$\star\star\star\star$}]
Un palet de masse \SI{0,160}{kg} glisse sur une glace horizontale.
Après une impulsion initiale, il se déplace en ligne droite à la vitesse constante \SI{8,0}{m.s^{-1}}.
Il percute ensuite une bande et repart sur la même droite en sens opposé à la vitesse
\SI{5,0}{m.s^{-1}}.

\begin{enumerate}[label=\arabic*.]
    \item Que peut-on dire de la résultante des forces extérieures pendant le glissement avant le choc ?
    \item Calculer la quantité de mouvement initiale du palet.
    \item Calculer sa quantité de mouvement finale après le choc.
    \item Déterminer la variation de quantité de mouvement.
    \item Interpréter le signe de cette variation.
    \item Expliquer qualitativement ce que cela dit de l’action de la bande sur le palet.
    \item Pourquoi cet exercice illustre-t-il bien la différence entre vitesse et quantité de mouvement ?
\end{enumerate}
\end{exercice}

\begin{exercice}[Type bac 10 -- Étude complète d’un tir de ballon \etoiles{$\star\star\star\star\star$}]
Un ballon est frappé depuis le sol avec une vitesse initiale de norme \SI{20}{m.s^{-1}} sous un angle
$\alpha=40^\circ$ avec l’horizontale. On néglige les frottements de l’air.

\begin{enumerate}[label=\arabic*.]
    \item Décomposer la vitesse initiale.
    \item Établir les équations horaires.
    \item Établir l’équation de la trajectoire.
    \item Déterminer le temps de montée.
    \item Déterminer l’altitude maximale atteinte.
    \item Déterminer la durée totale du vol.
    \item Déterminer la portée.
    \item Déterminer les composantes de la vitesse au sommet.
    \item Déterminer les composantes de la vitesse juste avant de retomber au sol.
    \item Montrer que la composante horizontale de la vitesse reste constante.
    \item Expliquer pourquoi la composante verticale s’annule au sommet alors que l’accélération ne s’annule pas.
    \item Discuter l’influence qualitative d’un angle de tir plus faible ou plus grand.
\end{enumerate}
\end{exercice}

%====================================================
\section{Exercices de synthèse niveau Terminale}
%====================================================

\begin{exercice}[Synthèse 1 -- Masse, poids, référentiel, forces \etoiles{$\star\star$}]
On s’intéresse à un skateur immobile sur une plate-forme, puis en mouvement rectiligne uniforme.

\begin{enumerate}[label=\arabic*.]
    \item Définir précisément les termes : système, référentiel, trajectoire, vitesse, accélération.
    \item Faire le bilan des forces lorsque le skateur est immobile.
    \item Que vaut la résultante des forces ?
    \item Même question lorsqu’il se déplace en ligne droite à vitesse constante.
    \item Expliquer pourquoi un mouvement est possible même si la somme des forces est nulle.
    \item Expliquer la différence entre “être en mouvement” et “être accéléré”.
\end{enumerate}
\end{exercice}

\begin{exercice}[Synthèse 2 -- Plan incliné sans frottement \etoiles{$\star\star\star$}]
Un objet de masse $m$ glisse sur un plan incliné d’angle $\alpha$, sans frottement.

\begin{enumerate}[label=\arabic*.]
    \item Faire un schéma clair du système et des forces.
    \item Décomposer le poids selon :
    \begin{enumerate}[label=\alph*)]
        \item l’axe parallèle au plan ;
        \item l’axe perpendiculaire au plan.
    \end{enumerate}
    \item Écrire la deuxième loi de Newton sur chacun des deux axes.
    \item En déduire l’expression de l’accélération.
    \item Comment varie cette accélération lorsque l’angle $\alpha$ augmente ?
    \item Que se passe-t-il pour $\alpha=0^\circ$ ? et pour $\alpha=90^\circ$ ?
\end{enumerate}
\end{exercice}

\begin{exercice}[Synthèse 3 -- Deux fils, un objet suspendu \etoiles{$\star\star\star$}]
Un projecteur de masse \SI{12}{kg} est suspendu immobile par deux câbles symétriques faisant chacun
un angle $\theta$ avec l’horizontale.

\begin{enumerate}[label=\arabic*.]
    \item Faire le bilan des forces.
    \item Pourquoi les tensions ont-elles même norme ?
    \item Décomposer chaque tension selon l’horizontale et la verticale.
    \item Écrire les conditions d’équilibre.
    \item En déduire une relation entre $T$, $m$, $g$ et $\theta$.
    \item Que se passe-t-il lorsque $\theta$ devient très petit ?
    \item Pourquoi cela explique-t-il qu’un câble presque horizontal peut être très tendu ?
\end{enumerate}
\end{exercice}

\begin{exercice}[Synthèse 4 -- Lecture d’un graphe $v(t)$ \etoiles{$\star\star\star$}]
On donne qualitativement un graphe vitesse-temps présentant trois phases :
\begin{itemize}
    \item de 0 à \SI{4}{s} : vitesse qui augmente linéairement ;
    \item de \SI{4}{s} à \SI{10}{s} : vitesse constante ;
    \item de \SI{10}{s} à \SI{14}{s} : vitesse qui diminue linéairement jusqu’à 0.
\end{itemize}

\begin{enumerate}[label=\arabic*.]
    \item Décrire le type de mouvement pendant chaque phase.
    \item Donner le signe de l’accélération pendant chaque phase.
    \item Dans quelle phase la résultante des forces est-elle nulle ?
    \item Comment reconnaître un mouvement uniformément accéléré sur un graphe $v(t)$ ?
    \item Comment reconnaître un mouvement uniforme ?
    \item Expliquer pourquoi une vitesse constante ne signifie pas forcément l’absence de forces dans toute situation réelle.
\end{enumerate}
\end{exercice}

\begin{exercice}[Synthèse 5 -- Action-réaction et propulsion \etoiles{$\star\star\star$}]
On étudie successivement un nageur, une fusée et une personne marchant sur le sol.

\begin{enumerate}[label=\arabic*.]
    \item Énoncer précisément la troisième loi de Newton.
    \item Dans le cas du nageur, identifier les deux forces d’action-réaction.
    \item Même question pour la marche.
    \item Même question pour la fusée.
    \item Expliquer pourquoi deux forces d’action-réaction ne s’annulent jamais pour un même système.
    \item Montrer que ce principe n’empêche absolument pas un objet d’accélérer.
\end{enumerate}
\end{exercice}

%====================================================
\section{Problèmes d’approfondissement}
%====================================================

\begin{exercice}[Approfondissement 1 -- Démonstration complète du lancer horizontal \etoiles{$\star\star\star\star$}]
On considère un projectile lancé horizontalement depuis une hauteur $h$ avec une vitesse initiale $v_0$.
On néglige les frottements.

\begin{enumerate}[label=\arabic*.]
    \item Faire le bilan des forces.
    \item Écrire la deuxième loi de Newton.
    \item Déterminer les coordonnées de l’accélération.
    \item Intégrer pour obtenir les composantes de la vitesse.
    \item Intégrer à nouveau pour obtenir les équations horaires.
    \item Montrer que le temps de chute vaut
    \[
    t_f=\sqrt{\frac{2h}{g}}.
    \]
    \item Montrer que la portée vaut
    \[
    L=v_0\sqrt{\frac{2h}{g}}.
    \]
    \item Montrer que le temps de chute ne dépend ni de la masse ni de la vitesse horizontale initiale.
    \item Déterminer les composantes de la vitesse à l’impact.
    \item Établir l’expression de la norme de la vitesse d’impact.
\end{enumerate}
\end{exercice}

\begin{exercice}[Approfondissement 2 -- Démonstration complète du lancer vertical \etoiles{$\star\star\star\star$}]
Un objet est lancé verticalement vers le haut avec une vitesse initiale $v_0$ depuis l’altitude $z_0$.

\begin{enumerate}[label=\arabic*.]
    \item Établir l’équation différentielle du mouvement.
    \item Déterminer $v(t)$ et $z(t)$.
    \item Déterminer l’instant du sommet.
    \item En déduire l’altitude maximale.
    \item Montrer que
    \[
    h_{\max}=\frac{v_0^2}{2g}
    \]
    si l’on mesure la hauteur par rapport au point de lancement.
    \item Expliquer pourquoi la vitesse s’annule au sommet mais pas l’accélération.
    \item Montrer que, sans frottements, la durée de montée est égale à la durée de descente jusqu’au niveau de départ.
\end{enumerate}
\end{exercice}

\begin{exercice}[Approfondissement 3 -- Mouvement circulaire uniforme et vitesse vectorielle \etoiles{$\star\star\star\star$}]
On considère un mobile décrivant un cercle de rayon $R$ à vitesse de norme constante $v$.

\begin{enumerate}[label=\arabic*.]
    \item Expliquer en quoi le mouvement est uniforme.
    \item Expliquer pourquoi le vecteur vitesse n’est pourtant pas constant.
    \item Décrire qualitativement l’orientation du vecteur vitesse en tout point.
    \item Décrire qualitativement l’orientation du vecteur accélération.
    \item Établir l’expression
    \[
    a_c=\frac{v^2}{R}.
    \]
    \item En déduire la résultante des forces nécessaire.
    \item Expliquer pourquoi une force perpendiculaire à la vitesse peut modifier la direction du mouvement
    sans modifier la norme de la vitesse.
\end{enumerate}
\end{exercice}

\begin{exercice}[Approfondissement 4 -- Quantité de mouvement et force résultante \etoiles{$\star\star\star\star$}]
Pour un système de masse constante, on note $\vec{p}=m\vec{v}$ sa quantité de mouvement.

\begin{enumerate}[label=\arabic*.]
    \item Montrer que
    \[
    \dv{\vec{p}}{t}=m\vec{a}.
    \]
    \item En déduire la forme
    \[
    \sum \vec{F}=\dv{\vec{p}}{t}.
    \]
    \item Interpréter physiquement le cas $\sum \vec{F}=\vec{0}$.
    \item En déduire une propriété sur la quantité de mouvement.
    \item Donner un exemple où la vitesse est non nulle et l’accélération nulle.
    \item Donner un exemple où la vitesse est nulle et l’accélération non nulle.
    \item Expliquer pourquoi ces deux situations ne sont pas contradictoires.
\end{enumerate}
\end{exercice}

\begin{exercice}[Approfondissement 5 -- Comparaison de modèles : avec et sans frottement \etoiles{$\star\star\star\star\star$}]
On étudie la chute d’une balle dans l’air.

\begin{enumerate}[label=\arabic*.]
    \item Dans le modèle de chute libre, quelles forces agissent sur la balle ?
    \item Quelle est alors l’accélération ?
    \item Ce modèle est-il plus pertinent pour une bille d’acier ou pour une feuille de papier ? Pourquoi ?
    \item Si l’on ajoute une force de frottement opposée au mouvement, que peut-on dire qualitativement :
    \begin{enumerate}[label=\alph*)]
        \item au début de la chute ;
        \item après un temps long ?
    \end{enumerate}
    \item Expliquer la notion de vitesse limite.
    \item Dire en quoi un modèle simple peut rester utile même s’il est imparfait.
\end{enumerate}
\end{exercice}

\begin{exercice}[Approfondissement 6 -- Étude complète d’un mouvement en deux phases \etoiles{$\star\star\star\star$}]
Un mobile de masse \SI{2,0}{kg} se déplace sur un plan horizontal sans frottement.
De $t=0$ à $t=\SI{3,0}{s}$, il subit une force horizontale constante de \SI{5,0}{N}.
À partir de $t=\SI{3,0}{s}$, cette force disparaît.

\begin{enumerate}[label=\arabic*.]
    \item Déterminer l’accélération pendant la première phase.
    \item Si le mobile part du repos, déterminer sa vitesse au bout de \SI{3,0}{s}.
    \item Déterminer la distance parcourue pendant la première phase.
    \item Décrire le mouvement pendant la seconde phase.
    \item Déterminer la vitesse pendant la seconde phase.
    \item Calculer la distance parcourue entre \SI{3,0}{s} et \SI{8,0}{s}.
    \item Représenter qualitativement l’allure des graphes $v(t)$ et $x(t)$.
\end{enumerate}
\end{exercice}

\begin{exercice}[Approfondissement 7 -- Problème de synthèse : virage puis départ en projectile \etoiles{$\star\star\star\star\star$}]
Un mobile de masse \SI{1,5}{kg} décrit d’abord un arc de cercle horizontal de rayon \SI{8,0}{m}
à la vitesse constante \SI{6,0}{m.s^{-1}}.
Il quitte ensuite la piste à l’horizontale depuis une hauteur \SI{2,5}{m}.

\begin{enumerate}[label=\arabic*.]
    \item Déterminer l’accélération centripète pendant la phase circulaire.
    \item Déterminer la résultante des forces pendant cette phase.
    \item Expliquer la direction de cette force.
    \item Lorsqu’il quitte la piste, quelle est sa vitesse initiale pour la phase de projectile ?
    \item Déterminer le temps de chute.
    \item Déterminer la portée horizontale.
    \item Déterminer les composantes de la vitesse à l’impact.
    \item Déterminer la norme de la vitesse à l’impact.
    \item Expliquer ce qui change et ce qui ne change pas entre la phase circulaire et la phase balistique.
\end{enumerate}
\end{exercice}

\begin{exercice}[Approfondissement 8 -- Problème ouvert de modélisation \etoiles{$\star\star\star\star\star$}]
On souhaite dimensionner grossièrement le lancer d’un ballon pour atteindre une cible située à
\SI{18}{m} du point de tir et à la même altitude.

On néglige l’air.

\begin{enumerate}[label=\arabic*.]
    \item Proposer un modèle physique raisonnable.
    \item Choisir les grandeurs pertinentes.
    \item Établir les équations du mouvement.
    \item Établir une relation entre la portée, la vitesse initiale et l’angle de tir.
    \item Discuter l’influence de l’angle.
    \item Expliquer pourquoi plusieurs couples $(v_0,\alpha)$ peuvent convenir.
    \item Discuter les limites physiques du modèle.
\end{enumerate}
\end{exercice}

\newpage

%====================================================
\section{Corrigé}
%====================================================

\subsection{Corrigé des Vrai / Faux}

\paragraph{VF 1}
\begin{enumerate}[label=\alph*)]
    \item Faux.
    \item Vrai.
    \item Vrai.
    \item Faux.
\end{enumerate}

\paragraph{VF 2}
\begin{enumerate}[label=\alph*)]
    \item Vrai.
    \item Faux.
    \item Vrai.
    \item Faux.
\end{enumerate}

\paragraph{VF 3}
\begin{enumerate}[label=\alph*)]
    \item Vrai.
    \item Faux.
    \item Vrai.
    \item Vrai.
\end{enumerate}

\paragraph{VF 4}
\begin{enumerate}[label=\alph*)]
    \item Vrai.
    \item Vrai.
    \item Faux.
    \item Vrai.
\end{enumerate}

\paragraph{VF 5}
\begin{enumerate}[label=\alph*)]
    \item Vrai.
    \item Faux.
    \item Faux.
    \item Vrai dans le référentiel d’étude habituel.
\end{enumerate}

\paragraph{VF 6}
\begin{enumerate}[label=\alph*)]
    \item Vrai.
    \item Vrai.
    \item Faux.
    \item Vrai.
\end{enumerate}

\paragraph{VF 7}
\begin{enumerate}[label=\alph*)]
    \item Vrai.
    \item Vrai.
    \item Faux.
    \item Vrai.
\end{enumerate}

\paragraph{VF 8}
\begin{enumerate}[label=\alph*)]
    \item Vrai.
    \item Vrai.
    \item Faux.
    \item Faux.
\end{enumerate}

\paragraph{VF 9}
\begin{enumerate}[label=\alph*)]
    \item Vrai.
    \item Faux.
    \item Vrai.
    \item Vrai.
\end{enumerate}

\paragraph{VF 10}
\begin{enumerate}[label=\alph*)]
    \item Vrai.
    \item Faux.
    \item Vrai.
    \item Vrai.
\end{enumerate}

\subsection{Corrigé des QCM}

\paragraph{QCM 1} C

\paragraph{QCM 2} B

\paragraph{QCM 3} C

\paragraph{QCM 4} B

\paragraph{QCM 5} C

\paragraph{QCM 6} C

\paragraph{QCM 7} B

\paragraph{QCM 8} C

\paragraph{QCM 9} C

\paragraph{QCM 10} B, C, D

\subsection{Corrigé des grands exercices type bac}

\paragraph{Type bac 1 -- Plongeon de haut vol}
\begin{enumerate}[label=\arabic*.]
    \item
    \[
    v_{0x}=v_0\cos\alpha=2,6\cos 35^\circ \approx \SI{2,13}{m.s^{-1}},
    \qquad
    v_{0y}=v_0\sin\alpha=2,6\sin 35^\circ \approx \SI{1,49}{m.s^{-1}}.
    \]
    \item Pendant la phase aérienne, on néglige l’air : seule la force poids agit.
    \[
    \vec P = m\vec g
    \]
    verticale vers le bas.
    \item Puisque seule la pesanteur agit, le plongeur est en chute libre.
    \item Deuxième loi de Newton :
    \[
    \sum \vec F = m\vec a \Rightarrow \vec P = m\vec a.
    \]
    Donc
    \[
    \vec a = \vec g = -g\,\vec j.
    \]
    Ainsi
    \[
    a_x=0,\qquad a_y=-g.
    \]
    \item En intégrant :
    \[
    v_x(t)=v_{0x},\qquad v_y(t)=v_{0y}-gt.
    \]
    \item Avec $x(0)=0$ et $y(0)=0$ :
    \[
    x(t)=v_{0x}t,\qquad y(t)=v_{0y}t-\frac12 gt^2.
    \]
    \item Comme $t=\dfrac{x}{v_{0x}}$,
    \[
    y(x)=\frac{v_{0y}}{v_{0x}}x-\frac{g}{2v_{0x}^2}x^2.
    \]
    \item C’est une fonction affine moins un terme en $x^2$ : la trajectoire est une parabole.
    \item L’eau est à l’altitude $y=-28$. On résout
    \[
    -28=v_{0y}t-\frac12 gt^2.
    \]
    Soit
    \[
    4,905t^2-1,49t-28=0.
    \]
    La solution positive vaut
    \[
    t \approx \SI{2,55}{s}.
    \]
    \item
    \[
    v_x \approx \SI{2,13}{m.s^{-1}},
    \qquad
    v_y \approx 1,49-9,81\times 2,55 \approx \SI{-23,5}{m.s^{-1}}.
    \]
    \item
    \[
    v=\sqrt{v_x^2+v_y^2}\approx \sqrt{2,13^2+23,5^2}\approx \SI{23,6}{m.s^{-1}}.
    \]
    \item L’affirmation “environ \SI{24}{m.s^{-1}}” est cohérente.
    \item La durée théorique vaut \SI{2,55}{s}. L’intervalle expérimental est
    \[
    [2,8-0,3;\,2,8+0,3]=[2,5;\,3,1]\ \si{s}.
    \]
    La valeur théorique appartient à cet intervalle : le modèle est compatible.
    \item Écarts possibles : frottements de l’air, posture du plongeur, vitesse initiale mal connue, hauteur réelle légèrement différente.
\end{enumerate}

\paragraph{Type bac 2 -- Service au volley-ball}
\begin{enumerate}[label=\arabic*.]
    \item
    \[
    v_{0x}=18\cos 12^\circ \approx \SI{17,6}{m.s^{-1}},
    \qquad
    v_{0y}=18\sin 12^\circ \approx \SI{3,74}{m.s^{-1}}.
    \]
    \item
    \[
    x(t)=v_{0x}t,\qquad y(t)=2,2+v_{0y}t-\frac12 gt^2.
    \]
    \item
    \[
    y(x)=2,2+\frac{v_{0y}}{v_{0x}}x-\frac{g}{2v_{0x}^2}x^2.
    \]
    \item À l’aplomb du filet, $x=9,0$ m :
    \[
    y(9)\approx 2,2+\frac{3,74}{17,6}\times 9-\frac{9,81}{2\times 17,6^2}\times 9^2.
    \]
    On trouve
    \[
    y(9)\approx \SI{2,83}{m}.
    \]
    \item Le filet mesure \SI{2,43}{m}, donc le ballon passe au-dessus.
    \item Pour le point de chute :
    \[
    0=2,2+3,74t-4,905t^2.
    \]
    La solution positive est
    \[
    t \approx \SI{1,16}{s}.
    \]
    Donc
    \[
    x_f=v_{0x}t \approx 17,6\times 1,16 \approx \SI{20,4}{m}.
    \]
    \item Le terrain adverse finit à \SI{18}{m}. Le ballon tombe au-delà : service dehors.
    \item Durée du vol : \SI{1,16}{s}.
    \item
    \[
    v_x=\SI{17,6}{m.s^{-1}},\qquad
    v_y=3,74-9,81\times1,16\approx \SI{-7,64}{m.s^{-1}}.
    \]
    \item
    \begin{enumerate}[label=\alph*)]
        \item Si $v_0$ augmente, la portée augmente en général.
        \item Si l’angle augmente légèrement, la hauteur augmente ; l’effet sur la portée dépend du cas.
        \item Si la hauteur initiale diminue, la durée de vol diminue donc la portée diminue.
    \end{enumerate}
    \item En réalité, il y a des frottements, éventuellement de la rotation, et la forme du ballon modifie la trajectoire.
\end{enumerate}

\paragraph{Type bac 3 -- Freinage d’urgence}
\begin{enumerate}[label=\arabic*.]
    \item
    \[
    90\,\si{km.h^{-1}} = \frac{90}{3,6}=\SI{25,0}{m.s^{-1}}.
    \]
    \item Distance de réaction :
    \[
    d_r=v_0t_r=25,0\times0,80=\SI{20,0}{m}.
    \]
    \item
    \[
    0=v_0+at \Rightarrow 0=25,0-6,5t
    \Rightarrow t=\frac{25,0}{6,5}\approx \SI{3,85}{s}.
    \]
    \item
    \[
    d_f=\frac{v_0^2}{2|a|}=\frac{25,0^2}{2\times 6,5}\approx \SI{48,1}{m}.
    \]
    \item
    \[
    d_{\text{arrêt}}=d_r+d_f\approx 20,0+48,1=\SI{68,1}{m}.
    \]
    \item
    \[
    \Delta \vec p = m(\vec v_f-\vec v_i)=1350(0-25,0)\vec i.
    \]
    Donc
    \[
    \Delta p = \SI{-33750}{kg.m.s^{-1}}.
    \]
    \item
    \[
    F_{\text{res}}=ma=1350\times(-6,5)=\SI{-8775}{N}.
    \]
    \item La résultante est opposée au mouvement.
    \item À \SI{110}{km.h^{-1}} :
    \[
    v_0=\frac{110}{3,6}\approx \SI{30,6}{m.s^{-1}}.
    \]
    Distance de réaction :
    \[
    d_r\approx 30,6\times0,80=\SI{24,4}{m}.
    \]
    Distance de freinage :
    \[
    d_f=\frac{30,6^2}{2\times6,5}\approx \SI{72,0}{m}.
    \]
    Distance totale :
    \[
    d_{\text{arrêt}}\approx \SI{96,4}{m}.
    \]
    \item L’augmentation est très importante : quelques km/h en plus augmentent beaucoup la distance d’arrêt.
    \item Car la distance de freinage dépend de $v_0^2$.
\end{enumerate}

\paragraph{Type bac 4 -- Virage sur route}
\begin{enumerate}[label=\arabic*.]
    \item
    \[
    54\,\si{km.h^{-1}}=\SI{15,0}{m.s^{-1}}.
    \]
    \item
    \[
    a_c=\frac{v^2}{R}=\frac{15^2}{45}=\SI{5,0}{m.s^{-2}}.
    \]
    \item
    \[
    F=ma=1000\times 5,0=\SI{5000}{N}.
    \]
    \item Cette résultante est horizontale, dirigée vers le centre du virage.
    \item La norme de la vitesse est constante, mais sa direction change : le vecteur vitesse varie.
    \item Si la vitesse double,
    \[
    a_c\propto v^2,
    \]
    donc elle est multipliée par 4.
    \item La force latérale aussi, car $F=ma_c$.
    \item Pour $R=90$ m :
    \[
    a_c=\frac{15^2}{90}=\SI{2,5}{m.s^{-2}},
    \]
    soit deux fois moins.
    \item Un virage plus serré demande une accélération centripète plus grande.
    \item Sur route mouillée, l’adhérence diminue, donc la force latérale maximale diminue.
\end{enumerate}

\paragraph{Type bac 5 -- Fusée au décollage}
\begin{enumerate}[label=\arabic*.]
    \item Forces : poussée vers le haut, poids vers le bas.
    \item
    \[
    P=mg=800\times 9,81=\SI{7848}{N}.
    \]
    \item
    \[
    F_{\text{res}}=12000-7848=\SI{4152}{N}.
    \]
    \item
    \[
    a=\frac{F_{\text{res}}}{m}=\frac{4152}{800}\approx \SI{5,19}{m.s^{-2}}.
    \]
    \item Si la poussée vaut \SI{7000}{N}, alors
    \[
    7000<7848,
    \]
    donc la résultante est vers le bas : pas de décollage.
    \item
    \[
    v=at=5,19\times 4,0\approx \SI{20,8}{m.s^{-1}}.
    \]
    \item
    \[
    z=\frac12 at^2=0,5\times5,19\times16\approx \SI{41,5}{m}.
    \]
    \item Les frottements réduiraient l’accélération.
    \item Car la fusée consomme du carburant : sa masse diminue avec le temps.
\end{enumerate}

\paragraph{Type bac 6 -- Particule chargée dans un champ électrique}
\begin{enumerate}[label=\arabic*.]
    \item Comme $q<0$, la force est opposée au champ.
    \item
    \[
    F_e=|q|E=5,0\times10^{-6}\times 800=\SI{4,0e-3}{N}.
    \]
    \item
    \[
    a=\frac{F}{m}=\frac{4,0\times10^{-3}}{2,0\times10^{-4}}=\SI{20}{m.s^{-2}}.
    \]
    \item C’est analogue à un projectile : accélération constante sur un axe perpendiculaire à la vitesse initiale.
    \item Durée de traversée :
    \[
    t=\frac{L}{v_0}=\frac{0,50}{12}\approx \SI{4,17e-2}{s}.
    \]
    \item Déviation verticale :
    \[
    y=\frac12 at^2=0,5\times20\times(4,17\times10^{-2})^2\approx \SI{1,74e-2}{m}.
    \]
    \item Poids :
    \[
    P=mg=2,0\times10^{-4}\times9,81\approx \SI{1,96e-3}{N}.
    \]
    La force électrique vaut \SI{4,0e-3}{N}, donc elle est environ deux fois plus grande.
    \item Le poids n’est pas totalement négligeable, mais l’ordre de grandeur reste inférieur à celui de la force électrique.
    \item L’hypothèse initiale est acceptable pour une première modélisation, mais pas parfaite.
\end{enumerate}

\paragraph{Type bac 7 -- Ascenseur}
\begin{enumerate}[label=\arabic*.]
    \item
    \[
    P=mg=70\times9,81\approx \SI{687}{N}.
    \]
    Ascenseur immobile : $R=P\approx \SI{687}{N}$.
    \item À vitesse constante, $a=0$, donc $R=P$.
    \item Accélération vers le haut :
    \[
    R-P=ma \Rightarrow R=m(g+a)=70(9,81+1,2)\approx \SI{771}{N}.
    \]
    \item Accélération vers le bas :
    \[
    R-P=-ma \Rightarrow R=m(g-a)=70(9,81-1,2)\approx \SI{603}{N}.
    \]
    \item La personne se sent plus lourde quand $R$ est plus grand : phase d’accélération vers le haut.
    \item Plus légère quand $R$ diminue : freinage en montée.
    \item Parce que la réaction traduit l’état dynamique vertical du système.
\end{enumerate}

\paragraph{Type bac 8 -- Chute Terre / Lune}
\begin{enumerate}[label=\arabic*.]
    \item
    \[
    h-\frac12 gt^2=0 \Rightarrow t=\sqrt{\frac{2h}{g}}.
    \]
    \item Sur Terre :
    \[
    t_T=\sqrt{\frac{20}{9,81}}\approx \SI{1,43}{s}.
    \]
    \item Sur la Lune :
    \[
    t_L=\sqrt{\frac{20}{1,62}}\approx \SI{3,51}{s}.
    \]
    \item
    \[
    v=\sqrt{2gh}.
    \]
    Donc
    \[
    v_T=\sqrt{2\times9,81\times10}\approx \SI{14,0}{m.s^{-1}},
    \qquad
    v_L=\sqrt{2\times1,62\times10}\approx \SI{5,69}{m.s^{-1}}.
    \]
    \item La masse ne change pas.
    \item Le poids est plus faible sur la Lune.
    \item Car c’est $g$ qui change, pas la masse.
\end{enumerate}

\paragraph{Type bac 9 -- Palet sur glace puis choc}
\begin{enumerate}[label=\arabic*.]
    \item Pendant le glissement uniforme sans frottement, la résultante des forces est nulle.
    \item
    \[
    p_i=mv_i=0,160\times 8,0=\SI{1,28}{kg.m.s^{-1}}.
    \]
    \item
    \[
    p_f=mv_f=0,160\times(-5,0)=\SI{-0,80}{kg.m.s^{-1}}.
    \]
    \item
    \[
    \Delta p=p_f-p_i=-0,80-1,28=\SI{-2,08}{kg.m.s^{-1}}.
    \]
    \item Le signe négatif indique une variation orientée en sens opposé au mouvement initial.
    \item La bande a exercé sur le palet une force impulsive importante opposée au mouvement initial.
    \item Parce que la quantité de mouvement tient compte à la fois de la masse, de la norme et du sens de la vitesse.
\end{enumerate}

\paragraph{Type bac 10 -- Étude complète d’un tir de ballon}
\begin{enumerate}[label=\arabic*.]
    \item
    \[
    v_{0x}=20\cos40^\circ\approx \SI{15,3}{m.s^{-1}},
    \qquad
    v_{0y}=20\sin40^\circ\approx \SI{12,9}{m.s^{-1}}.
    \]
    \item
    \[
    x(t)=v_{0x}t,\qquad y(t)=v_{0y}t-\frac12 gt^2.
    \]
    \item
    \[
    y(x)=\frac{v_{0y}}{v_{0x}}x-\frac{g}{2v_{0x}^2}x^2.
    \]
    \item Temps de montée :
    \[
    v_y=0 \Rightarrow t_m=\frac{v_{0y}}{g}\approx \frac{12,9}{9,81}\approx \SI{1,31}{s}.
    \]
    \item
    \[
    h_{\max}=\frac{v_{0y}^2}{2g}\approx \frac{12,9^2}{19,62}\approx \SI{8,48}{m}.
    \]
    \item Durée totale :
    \[
    T=2t_m\approx \SI{2,62}{s}.
    \]
    \item Portée :
    \[
    L=v_{0x}T\approx 15,3\times2,62\approx \SI{40,1}{m}.
    \]
    \item Au sommet :
    \[
    v_x=\SI{15,3}{m.s^{-1}},\qquad v_y=0.
    \]
    \item Juste avant de retomber :
    \[
    v_x=\SI{15,3}{m.s^{-1}},\qquad v_y\approx \SI{-12,9}{m.s^{-1}}.
    \]
    \item Car aucune force n’agit horizontalement si l’air est négligé, donc $a_x=0$.
    \item Au sommet, seule la composante verticale de la vitesse s’annule ; l’accélération reste $\vec g$ vers le bas.
    \item Un angle plus petit abaisse la trajectoire ; un angle plus grand augmente la hauteur mais pas toujours la portée.
\end{enumerate}

\subsection{Corrigé des exercices de synthèse niveau Terminale}

\paragraph{Synthèse 1}
\begin{enumerate}[label=\arabic*.]
    \item Système : objet étudié. Référentiel : solide de référence, repère et horloge.
    Trajectoire : ensemble des positions successives. Vitesse : variation de position.
    Accélération : variation du vecteur vitesse.
    \item Forces : poids et réaction du support.
    \item Résultante nulle.
    \item En MRU, la résultante est encore nulle.
    \item Parce qu’une vitesse constante est compatible avec une accélération nulle.
    \item Être en mouvement signifie avoir une vitesse non nulle ; être accéléré signifie que le vecteur vitesse varie.
\end{enumerate}

\paragraph{Synthèse 2}
\begin{enumerate}[label=\arabic*.]
    \item Forces : poids et réaction normale du plan.
    \item
    \[
    P_\parallel=mg\sin\alpha,\qquad P_\perp=mg\cos\alpha.
    \]
    \item Sur l’axe parallèle :
    \[
    ma=mg\sin\alpha.
    \]
    Sur l’axe perpendiculaire :
    \[
    R-mg\cos\alpha=0.
    \]
    \item
    \[
    a=g\sin\alpha.
    \]
    \item L’accélération augmente avec $\alpha$.
    \item Pour $\alpha=0^\circ$, $a=0$ ; pour $\alpha=90^\circ$, $a=g$.
\end{enumerate}

\paragraph{Synthèse 3}
\begin{enumerate}[label=\arabic*.]
    \item Poids et deux tensions.
    \item Par symétrie géométrique et équilibre.
    \item Les composantes horizontales se compensent ; les verticales s’additionnent.
    \item
    \[
    2T\sin\theta=mg
    \]
    si $\theta$ est l’angle avec l’horizontale.
    \item
    \[
    T=\frac{mg}{2\sin\theta}.
    \]
    \item Si $\theta$ devient très petit, $\sin\theta$ devient petit donc $T$ devient très grand.
    \item Un câble presque horizontal peut donc être très tendu même pour soutenir une masse modérée.
\end{enumerate}

\paragraph{Synthèse 4}
\begin{enumerate}[label=\arabic*.]
    \item Phase 1 : mouvement uniformément accéléré ; phase 2 : uniforme ; phase 3 : uniformément ralenti.
    \item Accélération positive, nulle, puis négative.
    \item Pendant la phase de vitesse constante.
    \item Par une droite oblique sur $v(t)$.
    \item Par une droite horizontale sur $v(t)$.
    \item En situation réelle, des forces peuvent exister mais se compenser.
\end{enumerate}

\paragraph{Synthèse 5}
\begin{enumerate}[label=\arabic*.]
    \item Deux systèmes exercent l’un sur l’autre des forces de même direction, même norme, sens opposés.
    \item Le nageur pousse l’eau vers l’arrière ; l’eau pousse le nageur vers l’avant.
    \item Le pied pousse le sol vers l’arrière ; le sol pousse la personne vers l’avant.
    \item La fusée éjecte les gaz vers le bas ; les gaz exercent une force vers le haut sur la fusée.
    \item Car elles agissent sur deux systèmes différents.
    \item Un objet accélère si la résultante des forces qui s’exercent sur lui n’est pas nulle.
\end{enumerate}

\subsection{Corrigé des problèmes d’approfondissement}

\paragraph{Approfondissement 1 -- Lancer horizontal}
\begin{enumerate}[label=\arabic*.]
    \item Seul le poids agit.
    \item
    \[
    m\vec a=\vec P=m\vec g.
    \]
    \item
    \[
    a_x=0,\qquad a_y=-g.
    \]
    \item
    \[
    v_x(t)=v_0,\qquad v_y(t)=-gt.
    \]
    \item
    \[
    x(t)=v_0t,\qquad y(t)=h-\frac12gt^2.
    \]
    \item
    \[
    0=h-\frac12gt_f^2
    \Rightarrow
    t_f=\sqrt{\frac{2h}{g}}.
    \]
    \item
    \[
    L=x(t_f)=v_0\sqrt{\frac{2h}{g}}.
    \]
    \item Le temps vient de l’équation verticale seulement.
    \item
    \[
    v_x=v_0,\qquad v_y=-gt_f=-\sqrt{2gh}.
    \]
    \item
    \[
    v=\sqrt{v_0^2+2gh}.
    \]
\end{enumerate}

\paragraph{Approfondissement 2 -- Lancer vertical}
\begin{enumerate}[label=\arabic*.]
    \item
    \[
    mz''(t)=-mg \Rightarrow z''(t)=-g.
    \]
    \item
    \[
    v(t)=v_0-gt,\qquad z(t)=z_0+v_0t-\frac12gt^2.
    \]
    \item Au sommet :
    \[
    v(t_s)=0 \Rightarrow t_s=\frac{v_0}{g}.
    \]
    \item
    \[
    z_{\max}=z_0+\frac{v_0^2}{2g}.
    \]
    \item Donc, par rapport au point de départ,
    \[
    h_{\max}=\frac{v_0^2}{2g}.
    \]
    \item Car l’accélération est due au poids, qui agit toujours.
    \item Par symétrie des équations en l’absence de frottements.
\end{enumerate}

\paragraph{Approfondissement 3 -- Mouvement circulaire uniforme}
\begin{enumerate}[label=\arabic*.]
    \item Uniforme car la norme de la vitesse est constante.
    \item Le vecteur vitesse change de direction.
    \item Il est tangent au cercle.
    \item L’accélération est dirigée vers le centre.
    \item
    \[
    a_c=\frac{v^2}{R}.
    \]
    \item
    \[
    F=ma_c=m\frac{v^2}{R}.
    \]
    \item Une force perpendiculaire change la direction de $\vec v$ sans modifier sa norme.
\end{enumerate}

\paragraph{Approfondissement 4 -- Quantité de mouvement}
\begin{enumerate}[label=\arabic*.]
    \item
    \[
    \vec p=m\vec v \Rightarrow \dv{\vec p}{t}=m\vec a
    \]
    si $m$ est constante.
    \item Donc
    \[
    \sum \vec F=\dv{\vec p}}{t}.
    \]
    \item Si $\sum\vec F=\vec 0$, alors $\vec p$ est constante.
    \item Donc la quantité de mouvement se conserve.
    \item Exemple : mouvement rectiligne uniforme.
    \item Exemple : sommet d’un lancer vertical.
    \item La vitesse peut être momentanément nulle sans que sa variation soit nulle.
\end{enumerate}

\paragraph{Approfondissement 5 -- Avec et sans frottement}
\begin{enumerate}[label=\arabic*.]
    \item Sans frottement : seul le poids agit.
    \item L’accélération vaut $\vec g$.
    \item Le modèle est meilleur pour une bille d’acier que pour une feuille.
    \item
    \begin{enumerate}[label=\alph*)]
        \item Au début, la vitesse augmente.
        \item Ensuite, on peut tendre vers une vitesse limite.
    \end{enumerate}
    \item La vitesse limite est atteinte quand le poids et le frottement se compensent.
    \item Un modèle simple permet déjà de comprendre les effets dominants.
\end{enumerate}

\paragraph{Approfondissement 6 -- Mouvement en deux phases}
\begin{enumerate}[label=\arabic*.]
    \item
    \[
    a=\frac{F}{m}=\frac{5,0}{2,0}=\SI{2,5}{m.s^{-2}}.
    \]
    \item
    \[
    v(3)=at=2,5\times 3,0=\SI{7,5}{m.s^{-1}}.
    \]
    \item
    \[
    x=\frac12 at^2=0,5\times2,5\times9=\SI{11,25}{m}.
    \]
    \item Ensuite, plus aucune force horizontale : mouvement rectiligne uniforme.
    \item Vitesse constante égale à \SI{7,5}{m.s^{-1}}.
    \item Entre 3 s et 8 s :
    \[
    d=7,5\times5=\SI{37,5}{m}.
    \]
    \item $v(t)$ : droite croissante puis plateau horizontal. $x(t)$ : parabole puis droite.
\end{enumerate}

\paragraph{Approfondissement 7 -- Virage puis projectile}
\begin{enumerate}[label=\arabic*.]
    \item
    \[
    a_c=\frac{6,0^2}{8,0}=\SI{4,5}{m.s^{-2}}.
    \]
    \item
    \[
    F=ma=1,5\times4,5=\SI{6,75}{N}.
    \]
    \item Vers le centre du cercle.
    \item La vitesse initiale de la phase projectile est horizontale de norme \SI{6,0}{m.s^{-1}}.
    \item
    \[
    t=\sqrt{\frac{2\times2,5}{9,81}}\approx \SI{0,714}{s}.
    \]
    \item
    \[
    L=6,0\times0,714\approx \SI{4,28}{m}.
    \]
    \item
    \[
    v_x=\SI{6,0}{m.s^{-1}},\qquad v_y=-gt\approx \SI{-7,00}{m.s^{-1}}.
    \]
    \item
    \[
    v=\sqrt{6,0^2+7,0^2}\approx \SI{9,22}{m.s^{-1}}.
    \]
    \item Dans la phase circulaire, il faut une force centripète ; dans la phase balistique, seule la pesanteur agit.
\end{enumerate}

\paragraph{Approfondissement 8 -- Problème ouvert de modélisation}
\begin{enumerate}[label=\arabic*.]
    \item On modélise le ballon comme un projectile ponctuel sans frottement.
    \item Grandeurs : $v_0$, $\alpha$, $g$, portée $L$.
    \item
    \[
    x(t)=v_0\cos\alpha\, t,\qquad
    y(t)=v_0\sin\alpha\, t-\frac12 gt^2.
    \]
    \item À altitude finale égale à l’altitude initiale, on obtient
    \[
    L=\frac{v_0^2}{g}\sin(2\alpha).
    \]
    \item L’angle agit via $\sin(2\alpha)$.
    \item Plusieurs couples $(v_0,\alpha)$ peuvent donner la même portée.
    \item Le modèle oublie l’air, la rotation, la taille du ballon et les imperfections du tir.
\end{enumerate}

\end{document}